\documentclass[11pt]{article}
\usepackage{mathblatt}
\def\mitzone{1}
\begin{document}
\blattkopf*{Prozentrechnung}{Gesamt}{Prozentrechnung · Gesamt · PRZ-L1}
\weit
\verzeichniszeile{\verz{zone}{Kennst du schon (Nr.~1–8)} \verztrenn \verz{e1}{1 Prozente als Anteile (Nr.~9–11)} \verztrenn \verz{e2}{2 Prozentsatz berechnen (Nr.~12–16)} \verztrenn \verz{e3}{3 Prozentwert berechnen (Nr.~17–19)} \verztrenn \verz{e4}{4 Grundwert berechnen (Nr.~20–22)} \verztrenn \verz{e5}{5 Prozentuale Veränderung (Nr.~23–25)} \verztrenn \verz{abhaken}{Das kann ich}}
% Zone „Das kennst du schon“ – aus bank/prozentrechnung/zone.jsonl (zusammenbau v0.3)
\einheitenkopf[zone][Kennst du schon]{Das kennst du schon}
\setcounter{aufgabe}{0}

%% TODO Titel: Ich-kann-Satz fehlt in der Bank (Platzhalter: Kettenname) – Nr. 1
%% TODO Anweisung: ein Satz über der ersten Teilaufgabe fehlt in der Bank – Nr. 1
\begin{aufgabe}{Bruch als Anteil lesen und kürzen (drei von zwölf, gekürzt ein Viertel), auf den Nenner hundert erweitern}
\begin{teile}
\teil $6$ von $10$ Bonbons sind rot – welcher Anteil? Kürze. \leerfeld
\teil $6$ von $15$ Plätzen sind besetzt – welcher Anteil? Kürze und erweitere auf Hundertstel. \leerfeld
\end{teile}
\end{aufgabe}

%% TODO Titel: Ich-kann-Satz fehlt in der Bank (Platzhalter: Kettenname) – Nr. 2
%% TODO Anweisung: ein Satz über der ersten Teilaufgabe fehlt in der Bank – Nr. 2
\begin{aufgabe}{Bruch ↔ Dezimalzahl (ein Viertel, drei Fünftel)}
\begin{teile}
\teil $\frac{1}{2}$ als Dezimalzahl? \leerfeld
\teil $\frac{3}{25}$ als Dezimalzahl? \leerfeld
\end{teile}
\end{aufgabe}

%% TODO Titel: Ich-kann-Satz fehlt in der Bank (Platzhalter: Kettenname) – Nr. 3
%% TODO Anweisung: ein Satz über der ersten Teilaufgabe fehlt in der Bank – Nr. 3
\begin{aufgabe}{Durch hundert teilen, mit Dezimalzahlen multiplizieren (Kommaverschiebung)}
\begin{teile}
\teil $600 : 100$ \leerfeld
\teil $0{,}15 \cdot 60$ \leerfeld
\end{teile}
\end{aufgabe}

%% TODO Titel: Ich-kann-Satz fehlt in der Bank (Platzhalter: Kettenname) – Nr. 4
%% TODO Anweisung: ein Satz über der ersten Teilaufgabe fehlt in der Bank – Nr. 4
\begin{aufgabe}{Hoch- und Runterrechnen in einer Tabelle (Dreisatz, proportionale Zuordnung)}
\begin{teile}
\teil $2$ Kinokarten kosten $18\,€$ – wie viel kosten $5$ Karten?

\begin{dreisatz}{Karten}{Preis}\dsz{2}{18\,€}\dsp{\phantom{:0}}{\phantom{:0}}\dsz{1}{\dsleer}\dsp{\phantom{:0}}{\phantom{:0}}\dsz{5}{\dsleer}\end{dreisatz}
\teil $5$ Brötchen kosten $2\,€$ – wie viel kosten $3$ Brötchen?

\begin{dreisatz}{Brötchen}{Preis}\dsz{5}{2\,€}\dsp{\phantom{:0}}{\phantom{:0}}\dsz{1}{\dsleer}\dsp{\phantom{:0}}{\phantom{:0}}\dsz{3}{\dsleer}\end{dreisatz}
\end{teile}
\end{aufgabe}

%% TODO Titel: Ich-kann-Satz fehlt in der Bank (Platzhalter: Kettenname) – Nr. 5
%% TODO Anweisung: ein Satz über der ersten Teilaufgabe fehlt in der Bank – Nr. 5
\begin{aufgabe}{Runden auf eine Dezimalstelle}
\begin{teile}
\teil $4{,}73$ – auf eine Dezimalstelle? \leerfeld
\teil $18{,}456$ kg – auf eine Dezimalstelle? \leerfeld[kg]
\end{teile}
\end{aufgabe}

%% TODO Titel: Ich-kann-Satz fehlt in der Bank (Platzhalter: Kettenname) – Nr. 6
%% TODO Anweisung: ein Satz über der ersten Teilaufgabe fehlt in der Bank – Nr. 6
\begin{aufgabe}{Bruchteil einer Größe (zwei Drittel von 60 €)}
\begin{teile}
\teil $\frac{1}{4}$ von $20\,€$ \leerfeld[€]
\teil $\frac{3}{5}$ von $45$ m \leerfeld[m]
\end{teile}
\end{aufgabe}

%% TODO Titel: Ich-kann-Satz fehlt in der Bank (Platzhalter: Kettenname) – Nr. 7
%% TODO Anweisung: ein Satz über der ersten Teilaufgabe fehlt in der Bank – Nr. 7
\begin{aufgabe}{Bruchteil einer Größe (zwei Drittel von 60 €) – Fehler finden}
\begin{teile}
\teil Von $40\,€$ ist ein Fünftel ausgegeben – wie viel ist übrig? Ben rechnet: $40 : 5 = 8$, übrig sind $8\,€$. Wo ist der Fehler?
\end{teile}
\end{aufgabe}

%% TODO Titel: Ich-kann-Satz fehlt in der Bank (Platzhalter: Kettenname) – Nr. 8
%% TODO Anweisung: ein Satz über der ersten Teilaufgabe fehlt in der Bank – Nr. 8
\begin{aufgabe}{Bruchteil einer Größe (zwei Drittel von 60 €) – selbst rechnen}
\begin{teile}
\teil Von $30$ Litern im Tank ist ein Drittel verbraucht – wie viel ist übrig? \leerfeld[l]
\end{teile}
\end{aufgabe}

\clearpage
% Einheit 1 von 5 – Prozente als Anteile (bank/prozentrechnung/e1.jsonl, zusammenbau v0.3)
%% TODO Zweigzeile Teil 1: was hier gelernt wird (Satz in Schülersprache aus dem Einheitstitel) – Einheit 1
\einheitenkopf[e1]{Einheit 1 von 5 · Prozente als Anteile}
\zweigzeile{ab Kl. 6 (am Gymnasium ab Kl. 5) · P10 oft}
\setcounter{aufgabe}{8}

%% TODO Titel: Ich-kann-Satz fehlt in der Bank (Platzhalter: Kettenname) – Nr. 9
%% TODO Anweisung: ein Satz über der ersten Teilaufgabe fehlt in der Bank – Nr. 9
\begin{aufgabe}{Umwandeln}
%% TODO \rechenplatz{n}: Schrittzahl des Grundfalls fehlt in der Bank (nur bei fester Schreibform, 2.3 b)
\begin{teile}
\teil Wie viel Prozent sind grau?

\streifenfeld[2]{50}{$0\,\%$}{$100\,\%$}
\teil $\frac{37}{100}$ – wie viel Prozent? \leerfeld[\%]
\teil $\frac{11}{20}$ – wie viel Prozent? \leerfeld[\%]
\teil $0{,}64$ – wie viel Prozent? \leerfeld[\%]
\teil Jeder hundertste Besucher bekommt einen Preis – wie viel Prozent? \leerfeld[\%]
\teil $35\,\%$ der Kinder fahren Rad, die anderen gehen zu Fuß – wie viel Prozent gehen zu Fuß? \leerfeld[\%]
\teil $5\,\%$ der Pakete kommen zu spät – welche Aussage passt? Kreuze an. \hfill (P10 2020 OS)\\ \kreuz{$5$ von $10$ Paketen kommen zu spät.}\\ \kreuz{Jedes 5. Paket kommt zu spät.}\\ \kreuz{$5$ von $100$ Paketen kommen zu spät.}
\teil Schulweg der Kinder einer Schule. Genau eine Aussage ist falsch. Kreuze sie an und schreibe sie richtig. \hfill (P10 2023 OS)\\ \kreuz{Aussage 1: Ein Viertel der Kinder fährt mit dem Bus.}\\ \kreuz{Aussage 2: Die Hälfte der Kinder fährt mit dem Rad.}

\sachtabelle{lr}{Schulweg & Kinder}{Bus & 150\\ Rad & 200\\ zu Fuß & 180\\ Auto & 70\\ gesamt & 600}
\teil Strom einer Gemeinde nach Energiequelle – trage den fehlenden Anteil ein. \hfill (P10 2021 OS)

\sachtabelle{lr}{Energiequelle & Anteil}{Wind & 38\,\%\\ Sonne & 24\,\%\\ Biogas & \leerzelle\\ Wasser & 9\,\%\\ Sonstiges & 3\,\%\\ gesamt & 100\,\%}
\end{teile}
\end{aufgabe}

%% TODO Titel: Ich-kann-Satz fehlt in der Bank (Platzhalter: Kettenname) – Nr. 10
%% TODO Anweisung: ein Satz über der ersten Teilaufgabe fehlt in der Bank – Nr. 10
\begin{aufgabe}{Prozentangaben ordnen}
\begin{teile}
\teil $0{,}4$; $\frac{1}{4}$; $35\,\%$; $\frac{3}{10}$ – der Größe nach, kleinste zuerst? \leerfeld $<$ \leerfeld $<$ \leerfeld $<$ \leerfeld
\end{teile}
\end{aufgabe}

%% TODO Titel: Ich-kann-Satz fehlt in der Bank (Platzhalter: Kettenname) – Nr. 11
%% TODO Anweisung: ein Satz über der ersten Teilaufgabe fehlt in der Bank – Nr. 11
\begin{aufgabe}{Umwandeln – Fehler finden, Begründen, Darstellungswechsel, Anwendung}
\begin{teile}
\teil Tom: „Jede fünfundzwanzigste Schraube ist fehlerhaft, das sind $25\,\%$.“ – Fehler? Wie heißt es richtig?
\teil Warum ist $\frac{1}{4}$ dasselbe wie $25\,\%$?
\teil Färbe $30\,\%$ des Streifens.

\streifenleer
\teil Abstimmung über den Wandertag: In der 7a sind $18$ von $25$ Kindern für den Zoo, in der 7b $21$ von $30$. In welcher Klasse ist der Anteil für den Zoo größer?
\end{teile}
\end{aufgabe}

\clearpage
% Einheit 2 von 5 – Prozentsatz berechnen (bank/prozentrechnung/e2.jsonl, zusammenbau v0.3)
%% TODO Zweigzeile Teil 1: was hier gelernt wird (Satz in Schülersprache aus dem Einheitstitel) – Einheit 2
\einheitenkopf[e2]{Einheit 2 von 5 · Prozentsatz berechnen}
\zweigzeile{ab Kl. 7 · P10 oft}
\setcounter{aufgabe}{11}

%% TODO Titel: Ich-kann-Satz fehlt in der Bank (Platzhalter: Kettenname) – Nr. 12
%% TODO Anweisung: ein Satz über der ersten Teilaufgabe fehlt in der Bank – Nr. 12
\begin{aufgabe}{Was ist das Ganze?}
\begin{teile}
\teil In der Klasse sind $28$ Kinder, $7$ davon haben einen Hund. – Was ist das Ganze?
\end{teile}
\end{aufgabe}

%% TODO Titel: Ich-kann-Satz fehlt in der Bank (Platzhalter: Kettenname) – Nr. 13
%% TODO Anweisung: ein Satz über der ersten Teilaufgabe fehlt in der Bank – Nr. 13
\begin{aufgabe}{Streifen einteilen}
\begin{teile}
\teil In $10$-\%-Schritte einteilen.

\streifenleer[0]
\end{teile}
\end{aufgabe}

%% TODO Titel: Ich-kann-Satz fehlt in der Bank (Platzhalter: Kettenname) – Nr. 14
%% TODO Anweisung: ein Satz über der ersten Teilaufgabe fehlt in der Bank – Nr. 14
\begin{aufgabe}{Prozentsatz}
%% TODO \rechenplatz{n}: Schrittzahl des Grundfalls fehlt in der Bank (nur bei fester Schreibform, 2.3 b)
\begin{teile}
\teil Streifen: $50$ Kinder, grau: $10$ Kinder – wie viel Prozent?

\streifenfeld{20}{$0$}{$50$ Kinder}
\teil $17$ von $100$ Schülern spielen ein Instrument – wie viel Prozent? \leerfeld[\%]
\teil $19$ von $50$ Kindern – wie viel Prozent? \leerfeld[\%]
\teil $27$ von $60$ Minuten Training sind Laufen – wie viel Prozent? \leerfeld[\%]
\teil $17$ von $24$ Kindern waren im Schwimmbad – wie viel Prozent? Runde auf eine Stelle. \leerfeld[\%]
\teil Im Chor singen $14$ Mädchen und $6$ Jungen – wie viel Prozent sind Jungen? \leerfeld[\%]
\teil Jacke: vorher $80\,€$, jetzt $60\,€$ – wie viel Prozent Rabatt? \leerfeld[\%]
\teil Bei einer Tombola gibt es $21$ Nieten und $3$ Hauptgewinne. Wie viel Prozent der Lose sind Hauptgewinne? \hfill (P10 2018 OS) \\ \leerfeld[\%]
\teil Eine Stadt hat $40\,000$ Einwohner, $1\,080$ sind in diesem Jahr zugezogen. Wie viel Prozent sind das? \hfill (P10 2014 OS) \\ \leerfeld[\%]
\teil Körpergrößen der $13$ Spieler einer Mannschaft in cm: $172$, $185$, $169$, $190$, $178$, $181$, $166$, $188$, $175$, $183$, $170$, $192$, $177$. Wie viel Prozent sind größer als $180$ cm? Runde auf eine Stelle. \hfill (P10 2025 OS)
\end{teile}
\end{aufgabe}

%% TODO Titel: Ich-kann-Satz fehlt in der Bank (Platzhalter: Kettenname) – Nr. 15
%% TODO Anweisung: ein Satz über der ersten Teilaufgabe fehlt in der Bank – Nr. 15
\begin{aufgabe}{Umkehrung: zu einem Prozentsatz ein Zahlenpaar angeben}
\begin{teile}
\teil $30\,\%$ – Teil und Ganzes? \leerfeld von \leerfeld
\end{teile}
\end{aufgabe}

%% TODO Titel: Ich-kann-Satz fehlt in der Bank (Platzhalter: Kettenname) – Nr. 16
%% TODO Anweisung: ein Satz über der ersten Teilaufgabe fehlt in der Bank – Nr. 16
\begin{aufgabe}{Prozentsatz – Fehler finden, Begründen, Darstellungswechsel, Anwendung}
\begin{teile}
\teil Kim, „$12$ von $48$ Kindern“: \rechnung{48 : 12 &= 4 = 400\,\%} Fehler? Wie heißt es richtig?
\teil Warum teilt man beim Prozentsatz durch das Ganze und nicht durch den Teil?
\teil Ganzes $40$ kg, Teil $10$ kg – im Streifen einzeichnen, Prozentsatz ablesen.

\streifenleer[4]
\leerfeld[\%]
\teil Handyvertrag mit $20$ GB im Monat. Am 20. Tag sind $17$ GB verbraucht. Wie viel Prozent sind noch übrig? Reicht der Rest bei gleichem Verbrauch für die letzten $10$ Tage?
\end{teile}
\end{aufgabe}

\clearpage
% Einheit 3 von 5 – Prozentwert berechnen (bank/prozentrechnung/e3.jsonl, zusammenbau v0.3)
%% TODO Zweigzeile Teil 1: was hier gelernt wird (Satz in Schülersprache aus dem Einheitstitel) – Einheit 3
\einheitenkopf[e3]{Einheit 3 von 5 · Prozentwert berechnen}
\zweigzeile{ab Kl. 7 · P10}
\setcounter{aufgabe}{16}

%% TODO Titel: Ich-kann-Satz fehlt in der Bank (Platzhalter: Kettenname) – Nr. 17
%% TODO Anweisung: ein Satz über der ersten Teilaufgabe fehlt in der Bank – Nr. 17
\begin{aufgabe}{Prozentwert}
%% TODO \rechenplatz{n}: Schrittzahl des Grundfalls fehlt in der Bank (nur bei fester Schreibform, 2.3 b)
\begin{teile}
\teil $50\,\%$ von $80\,€$ – wie viel?

\streifen[2]{50}{$0$}{$80\,€$}
\leerfeld[€]
\teil $6\,\%$ von $400$ kg \leerfeld[kg]
\teil $30\,\%$ von $90\,€$ \leerfeld[€]
\teil $45\,\%$ von $80$ kg – mit Dezimalzahl \leerfeld[kg]
\teil $20\,\%$ von $4{,}50\,€$ \leerfeld[€]
\teil Schuhe für $70\,€$, $20\,\%$ Rabatt – wie viel spart man? \leerfeld[€]
\teil $150\,\%$ von $40$ kg \leerfeld[kg]
\teil Ein Fernseher kostet $480\,€$. Bei Barzahlung gibt es $15\,\%$ Rabatt. Wie viel Euro spart man? Kreuze an. \hfill (P10 2021 OS)\\ \kreuz{$32\,€$}\\ \kreuz{$408\,€$}\\ \kreuz{$72\,€$}\\ \kreuz{$552\,€$}
\teil $17\,\%$ von $60\,€$ \hfill (P10 2014 OS) \\ \leerfeld[€]
\teil Bei einer Umfrage unter $800$ Haushalten hatten $35\,\%$ ein E-Bike, sechs Jahre später $52\,\%$. Um wie viele Haushalte hat die Zahl zugenommen? \hfill (P10 2019 OS)
\end{teile}
\end{aufgabe}

%% TODO Titel: Ich-kann-Satz fehlt in der Bank (Platzhalter: Kettenname) – Nr. 18
%% TODO Anweisung: ein Satz über der ersten Teilaufgabe fehlt in der Bank – Nr. 18
\begin{aufgabe}{Mehrwertsteuer in Euro}
\begin{teile}
\teil Reparatur ohne Mehrwertsteuer $300\,€$, dazu $19\,\%$ – wie viel Euro Steuer? \leerfeld[€]
\end{teile}
\end{aufgabe}

%% TODO Titel: Ich-kann-Satz fehlt in der Bank (Platzhalter: Kettenname) – Nr. 19
%% TODO Anweisung: ein Satz über der ersten Teilaufgabe fehlt in der Bank – Nr. 19
\begin{aufgabe}{Prozentwert – Fehler finden, Begründen, Darstellungswechsel, Anwendung}
\begin{teile}
\teil Jonas, $60\,\%$ von $35\,€$: \rechnung{0{,}06 \cdot 35 &= 2{,}10\,€} Fehler? Wie heißt es richtig?
\teil $1\,\%$ von $500\,€$ sind $5\,€$. Warum?
\teil Ganzes $60\,€$: $30\,\%$ im Streifen einzeichnen, Wert ablesen.

\streifen[0]{0}{$0$}{$60\,€$}
\leerfeld[€]
\teil Dieselbe Jacke: in Laden A $95\,€$ mit $20\,\%$ Rabatt, in Laden B $85\,€$ mit $10\,\%$ Rabatt. Wo ist sie billiger?
\end{teile}
\end{aufgabe}

\clearpage
% Einheit 4 von 5 – Grundwert berechnen (bank/prozentrechnung/e4.jsonl, zusammenbau v0.3)
%% TODO Zweigzeile Teil 1: was hier gelernt wird (Satz in Schülersprache aus dem Einheitstitel) – Einheit 4
\einheitenkopf[e4]{Einheit 4 von 5 · Grundwert berechnen}
\zweigzeile{ab Kl. 7 · P10 oft}
\setcounter{aufgabe}{19}

%% TODO Titel: Ich-kann-Satz fehlt in der Bank (Platzhalter: Kettenname) – Nr. 20
%% TODO Anweisung: ein Satz über der ersten Teilaufgabe fehlt in der Bank – Nr. 20
\begin{aufgabe}{Was ist gegeben, was gesucht?}
\begin{teile}
\teil $30\,\%$ der Klasse sind $9$ Kinder. Wie viele Kinder hat die Klasse? – Was ist gesucht?\\ \kreuz{Prozentwert}\\ \kreuz{Prozentsatz}\\ \kreuz{Grundwert}
\end{teile}
\end{aufgabe}

%% TODO Titel: Ich-kann-Satz fehlt in der Bank (Platzhalter: Kettenname) – Nr. 21
%% TODO Anweisung: ein Satz über der ersten Teilaufgabe fehlt in der Bank – Nr. 21
\begin{aufgabe}{Grundwert}
%% TODO \rechenplatz{n}: Schrittzahl des Grundfalls fehlt in der Bank (nur bei fester Schreibform, 2.3 b)
\begin{teile}
\teil Grau sind $30\,\%$, das sind $6$ Kästchen – wie viele Kästchen hat der ganze Streifen?

\streifen[0]{30}{$0\,\%$}{$100\,\%$}
\leerfeld[Kästchen]
\teil $50\,\%$ sind $35$ kg – wie viel ist das Ganze? \leerfeld[kg]
\teil $3\,\%$ sind $12$ kg – wie viel ist das Ganze? \leerfeld[kg]
\teil $36\,\%$ sind $81$ kg – wie viel ist das Ganze? \leerfeld[kg]
\teil $21$ Kinder fahren mit dem Rad, das sind $75\,\%$ der Klasse – wie viele Kinder hat die Klasse? \leerfeld
\teil $40\,\%$ von $65$ Kindern gehen in die AG – wie viele? \leerfeld
\teil Beim Kauf einer Mütze spart Mira $4\,€$, das sind $25\,\%$ des alten Preises. Wie viel kostete die Mütze vorher? \hfill (P10 2025 OS) \\ \leerfeld[€]
\end{teile}
\end{aufgabe}

%% TODO Titel: Ich-kann-Satz fehlt in der Bank (Platzhalter: Kettenname) – Nr. 22
%% TODO Anweisung: ein Satz über der ersten Teilaufgabe fehlt in der Bank – Nr. 22
\begin{aufgabe}{Grundwert – Fehler finden, Begründen, Darstellungswechsel, Anwendung}
\begin{teile}
\teil „$30\,\%$ sind $12$ kg – wie viel ist das Ganze?“ Sven: \rechnung{0{,}3 \cdot 12 &= 3{,}6\text{ kg}} Fehler? Wie heißt es richtig?
\teil Warum rechnet man beim Grundwert am Ende mal $100$?
\teil Grau sind $40\,\%$, das sind $18$ kg. Trage die Werte an den Streifen, lies das Ganze ab.

\streifen{40}{$0$ kg}{?}
\leerfeld[kg]
\teil Tims Handy zeigt $35\,\%$ Akku; laut Anzeige reicht das noch für $7$ Stunden Musik. Wie lange hält ein voller Akku? Reicht er für eine Busfahrt von $18$ Stunden?
\end{teile}
\end{aufgabe}

\clearpage
% Einheit 5 von 5 – Prozentuale Veränderung (bank/prozentrechnung/e5.jsonl, zusammenbau v0.3)
%% TODO Zweigzeile Teil 1: was hier gelernt wird (Satz in Schülersprache aus dem Einheitstitel) – Einheit 5
\einheitenkopf[e5]{Einheit 5 von 5 · Prozentuale Veränderung}
\zweigzeile{ab Kl. 7, je nach Buch bis Kl. 8 · P10 oft}
\setcounter{aufgabe}{22}

%% TODO Titel: Ich-kann-Satz fehlt in der Bank (Platzhalter: Kettenname) – Nr. 23
%% TODO Anweisung: ein Satz über der ersten Teilaufgabe fehlt in der Bank – Nr. 23
\begin{aufgabe}{Welcher Wert ist der alte?}
\begin{teile}
\teil Der Preis stieg von $2{,}40\,€$ auf $2{,}70\,€$. – Welcher Wert ist der alte?
\end{teile}
\end{aufgabe}

%% TODO Titel: Ich-kann-Satz fehlt in der Bank (Platzhalter: Kettenname) – Nr. 24
%% TODO Anweisung: ein Satz über der ersten Teilaufgabe fehlt in der Bank – Nr. 24
\begin{aufgabe}{Veränderung}
%% TODO \rechenplatz{n}: Schrittzahl des Grundfalls fehlt in der Bank (nur bei fester Schreibform, 2.3 b)
\begin{teile}
\teil um $20\,\%$ gesenkt – was bedeutet dasselbe? Kreuze an.\\ \kreuz{auf $20\,\%$ gesenkt}\\ \kreuz{auf $80\,\%$ gesenkt}\\ \kreuz{um $80\,\%$ gesenkt}
\teil $80\,€$ um $10\,\%$ erhöht – neuer Preis? \leerfeld[€]
\teil $45\,€$ um $20\,\%$ erhöht – mit Faktor \leerfeld[€]
\teil von $40$ auf $50$ – um wie viel Prozent mehr? \leerfeld[\%]
\teil Nach $20\,\%$ Rabatt kostet ein Rad $360\,€$ – alter Preis? \leerfeld[€]
\teil netto (ohne Steuer) $200\,€$, $19\,\%$ Mehrwertsteuer – brutto? \leerfeld[€]
\teil Wahlbeteiligung $58\,\%$, später $64\,\%$ – wie viele Prozentpunkte mehr? Wie viel Prozent mehr? Runde auf eine Stelle.
\teil Eine Straße steigt auf $200$ m waagerechter Strecke um $14$ m – Steigung in Prozent? \leerfeld[\%]
\teil Ein Liter Milch kostete $1{,}15\,€$, jetzt kostet er $1{,}29\,€$. Um wie viel Prozent ist der Preis gestiegen? Runde auf eine Stelle. \hfill (P10 2026 FOR)
\teil Ein Busticket kostet $2{,}40\,€$ und wird um $25\,\%$ teurer. Wie viel kostet es dann? \hfill (P10 2024 OS) \\ \leerfeld[€]
\teil Kletterhalle: Erwachsene $14\,€$, Jugendliche (bis $17$ Jahre) $9\,€$. Familie Demir: Mutter, Vater, Tochter ($15$ Jahre), Sohn ($12$ Jahre). Montags gibt es $15\,\%$ Rabatt. Wie viel zahlt die Familie am Montag? \hfill (P10 2015 OS)
\teil Für Radwege gilt: Steigung höchstens $5\,\%$. Ergänze: „$5\,\%$ Steigung heißt: auf \leerfeld[m] waagerechter Strecke \leerfeld[m] Höhenunterschied.“ Herr Lenz sagt: „Ein Radweg, der auf $240$ m waagerechter Strecke um $15$ m ansteigt, ist zu steil.“ Hat er recht? Rechne. \hfill (P10 2025 OS)
\end{teile}
\end{aufgabe}

%% TODO Titel: Ich-kann-Satz fehlt in der Bank (Platzhalter: Kettenname) – Nr. 25
%% TODO Anweisung: ein Satz über der ersten Teilaufgabe fehlt in der Bank – Nr. 25
\begin{aufgabe}{Veränderung – Fehler finden, Begründen, Anwendung}
\begin{teile}
\teil Preis von $60\,€$ auf $75\,€$ – Steigerung in Prozent? Ida: \rechnung{15 : 75 &= 0{,}2 = 20\,\%} Fehler? Wie heißt es richtig?
\teil Warum bezieht man eine Preiserhöhung auf den alten Preis?
\teil Zwei Angebote für ein Fahrrad zu $600\,€$: A gibt $10\,\%$ Rabatt und auf den neuen Preis nochmals $10\,\%$; B gibt einmal $20\,\%$ Rabatt. Welches Angebot ist günstiger?
\end{teile}
\end{aufgabe}

% Abhakseite „Das kann ich“ – Zone nur im Gesamt (\mitzone)
%% TODO Ich-kann-Sätze: Platzhalter wie in den Aufgabendateien
\begin{abhakseite}
\ifdefined\mitzone
\abhakgruppe{Kennst du schon}
\abhak{1}{Bruch als Anteil lesen und kürzen (drei von zwölf, gekürzt ein Viertel), auf den Nenner hundert erweitern}
\abhak{2}{Bruch ↔ Dezimalzahl (ein Viertel, drei Fünftel)}
\abhak{3}{Durch hundert teilen, mit Dezimalzahlen multiplizieren (Kommaverschiebung)}
\abhak{4}{Hoch- und Runterrechnen in einer Tabelle (Dreisatz, proportionale Zuordnung)}
\abhak{5}{Runden auf eine Dezimalstelle}
\abhak{6}{Bruchteil einer Größe (zwei Drittel von 60 €)}
\abhak{7}{Bruchteil einer Größe (zwei Drittel von 60 €) – Fehler finden}
\abhak{8}{Bruchteil einer Größe (zwei Drittel von 60 €) – selbst rechnen}
\fi
\abhakgruppe{Einheit 1 · Prozente als Anteile}
\abhak{9}{Umwandeln}
\abhak{10}{Prozentangaben ordnen}
\abhak{11}{Umwandeln – Fehler finden, Begründen, Darstellungswechsel, Anwendung}
\abhakgruppe{Einheit 2 · Prozentsatz berechnen}
\abhak{12}{Was ist das Ganze?}
\abhak{13}{Streifen einteilen}
\abhak{14}{Prozentsatz}
\abhak{15}{Umkehrung: zu einem Prozentsatz ein Zahlenpaar angeben}
\abhak{16}{Prozentsatz – Fehler finden, Begründen, Darstellungswechsel, Anwendung}
\abhakgruppe{Einheit 3 · Prozentwert berechnen}
\abhak{17}{Prozentwert}
\abhak{18}{Mehrwertsteuer in Euro}
\abhak{19}{Prozentwert – Fehler finden, Begründen, Darstellungswechsel, Anwendung}
\abhakgruppe{Einheit 4 · Grundwert berechnen}
\abhak{20}{Was ist gegeben, was gesucht?}
\abhak{21}{Grundwert}
\abhak{22}{Grundwert – Fehler finden, Begründen, Darstellungswechsel, Anwendung}
\abhakgruppe{Einheit 5 · Prozentuale Veränderung}
\abhak{23}{Welcher Wert ist der alte?}
\abhak{24}{Veränderung}
\abhak{25}{Veränderung – Fehler finden, Begründen, Anwendung}
\end{abhakseite}

\end{document}
